\section{Conclusion}\label{sec:conclusion}

Every evaluation in \ours{}, including its own headline number, treats the
verifier as a judge: score $k$ draws, report the best. Feeding the verifier's
diagnosis back into regeneration, on an untuned base \texttt{gpt-oss-120b} and
at a lower model-call budget than the open-loop control, raises the mean
solved count on the frozen 30-spec holdout from 10.7/30 to 17.3/30 across
three seeds each, with non-overlapping ranges. It also exceeds 15/30, the best
open-loop result this training program has produced with fine-tuning, a
result that the program's own pre-registered rule already marks null. Half of
framing L's solves in the seed with per-call provenance come from a fresh
draw and half from a repair round; five of the six specs gained over the
open-loop control were won on repair, not on generation. Only two of the 30
holdout specs are solved by every closed-loop seed and no open-loop seed; the
rest of the gain comes from a larger reachable set whose composition shifts
by seed, not from a fixed set of specs the loop alone can reach.

The failure taxonomy and the decode-time grammar bound what closing the loop
is working against and what remains unaddressed. 51\% of pooled SANY failures
do not parse, and a grammar built for the \tlaplus{} expression language
rejects 86.7\% of 1,543 real parse failures with zero false rejects on 438
known-good specs. That number describes what a decode-time constraint would
remove from the failure population, not a measured pass-rate gain: the
grammar has not been run inside a generation arm, and a companion lint pass
that forces 88 failing candidates to parse converts zero of them to a TLC
pass. Forcing a candidate past the parser does not by itself get it past the
model checker. The 2x2 this addendum motivates, framing (open vs. closed) by
training (base vs. fine-tuned), has one cell unmeasured: framing L has run on
the base model only, and the tuned-model cell requires a self-hosted serve
that has not yet run, so nothing here says what task training adds on top of
the loop. A related but separate gap is decode-time constraint itself: the
shared inference endpoint used for every arm in this addendum accepts vLLM's
\texttt{guided\_grammar} and \texttt{guided\_choice} parameters with HTTP 200
and applies neither, so the grammar's 86.7\% catch rate cannot yet be turned
into a generation-arm result on this endpoint. Closing either gap requires a
self-hosted serve, which has not yet run.
